% Zdroj: PDF priklady-podzim-2022.pdf, fyzická str. 1-2. Značka: společné okruhy. Zařazeno: I4. Číselné výsledky ověřeny numericky.
\section{Virtuální paměť}
\szzotazka{podzim 2022}{2}{Virtuální paměť (stránkování, 32-bit)}

\noindent\textbf{Zadání.}\par
Uvažujme procesor s~podporou stránkování. Virtuální i~fyzická adresa má 32 bitů, používá
se jednoúrovňová stránkovací tabulka a~velikost stránky je 4\,KiB. Vedle nezbytných atributů
lze pro každou stránku nastavit, zda ji smí proces modifikovat a~zda v~ní lze spouštět kód.
Procesor navíc poskytuje informaci, zda se ke stránce přistupovalo a~zda do ní bylo zapsáno.

\medskip
\noindent\textbf{Část A.}
\begin{enumerate}
\item Navrhněte a~schematicky znázorněte formát položky stránkovací tabulky a~vysvětlete
  význam jednotlivých polí. Položka musí obsahovat všechny nutné informace a~být efektivně
  přístupná (nejvýše 1 čtení z~paměti).
\item Kolik položek bude mít stránkovací tabulka a~kolik paměti zabere pro každý proces?
\end{enumerate}

\noindent\textbf{Část B.}
\begin{enumerate}
\item Napište (v~pseudokódu) funkci \texttt{create\_pte} se signaturou
  \texttt{pte\_t create\_pte(address\_t phys\_addr, bool present, bool writable, bool executable)},
  která pro zadanou fyzickou adresu a~atributy vrátí položku stránkovací tabulky.
\item Napište (v~pseudokódu) funkci \texttt{void set\_mapping(pte\_t[] page\_table, address\_t virt\_addr, pte\_t entry)},
  která do tabulky pro zadanou virtuální adresu vloží záznam.
\end{enumerate}

\noindent\textbf{Část C.}
Souvislý blok 4 stránek virtuální paměti od adresy \texttt{0x0000B000} byl namapován na
souvislý blok fyzické paměti. Blok je část heapu (lze číst a~modifikovat, nelze v~něm
spouštět kód). Proces přečetl proměnnou (4\,B) z~virtuální adresy \texttt{0x0000CE90}, což
bylo přeloženo na fyzickou adresu \texttt{0xA1018E90}. Pozměněnou hodnotu zapsal na
virtuální adresu \texttt{0x0000D854}.
\begin{enumerate}
\item Na kterých fyzických adresách jsou uloženy položky stránkovací tabulky pro tento blok,
  je-li tabulka v~souvislém bloku fyzické paměti od adresy \texttt{0x13400000}?
\item Uveďte hexadecimální hodnoty všech položek tabulky pro tento rozsah, včetně všech atributů.
\end{enumerate}

V~odpovědi důsledně uvádějte číselnou soustavu -- prefix \texttt{0x} označuje šestnáctkovou soustavu.

\medskip
\noindent\textbf{Náčrt řešení.}\par

\noindent\textit{Část A.2 (rozměry tabulky).}
Stránka má $4\,\text{KiB}=2^{12}$\,B, takže offset uvnitř stránky tvoří dolních
$\boxed{12}$ bitů a~číslo stránky horních $32-12=\boxed{20}$ bitů. Stránkovací tabulka má
proto $2^{20}=\boxed{1\,048\,576}$ položek. Při velikosti položky 4\,B (32 bitů) zabere
tabulka $2^{20}\cdot 4\,\text{B}=\boxed{4\,\text{MiB}}$ na proces.

\medskip
\noindent\textit{Část A.1 (návrh položky -- PTE, 32 bitů).}
Číslo fyzického rámce (PFN) má rovněž $32-12=20$ bitů, vejde se tedy do horních 20 bitů
položky [31:12]. Dolních 12 bitů (které by u~adresy tvořily offset) zbývá na příznaky --
těch potřebujeme jen 5, takže 4-bytová položka bohatě stačí a~celá se přečte jedním přístupem
do paměti. Rozložení polí je volbou návrháře; v~dalším pevně zafixujeme toto:
\begin{center}
\begin{tabular}{|c|c|c|c|c|c|}
\hline
\textbf{bity [31:12]} & \textbf{bit 6} & \textbf{bit 5} & \textbf{bit 2} & \textbf{bit 1} & \textbf{bit 0} \\
\hline
PFN (číslo rámce) & D & A & X & W & P \\
\hline
\end{tabular}
\end{center}
kde \textbf{P} = Present (platnost záznamu), \textbf{W} = Writable (proces smí stránku
modifikovat), \textbf{X} = eXecutable (lze v~ní spouštět kód), \textbf{A} = Accessed
(ke stránce se přistupovalo), \textbf{D} = Dirty (do stránky bylo zapsáno). Bity 3, 4, 7--11
zůstávají rezervované (zde 0). Bity A a~D nastavuje hardware, ostatní spravuje OS.

\medskip
\noindent\textit{Část B (pseudokód).}
\begin{lstlisting}
function create_pte(phys_addr, present, writable, executable):
    // PFN je v hornich 20 bitech, dolni offsetove bity vynulujeme
    pte = phys_addr & 0xFFFFF000
    // priznaky vskladame podle zvoleneho rozlozeni z casti A
    if present:    pte = pte | 0x1   // bit 0: Present
    if writable:   pte = pte | 0x2   // bit 1: Writable
    if executable: pte = pte | 0x4   // bit 2: eXecutable
    return pte

procedure set_mapping(page_table, virt_addr, entry):
    // cislo stranky = hornich 20 bitu virtualni adresy
    index = virt_addr >> 12
    page_table[index] = entry
\end{lstlisting}

\medskip
\noindent\textit{Část C (překlad a~položky tabulky).}
Stránka má \texttt{0x1000}\,B, číslo stránky je horních 20 bitů adresy (tj. adresa
\texttt{>> 12}), offset dolních 12 bitů.

Čtená virtuální adresa \texttt{0x0000CE90}: číslo stránky $=\texttt{0xC}$, offset
$=\texttt{0xE90}$. Fyzická adresa \texttt{0xA1018E90} má stejný offset \texttt{0xE90}
(souhlasí) a~rámec $=\texttt{0xA1018}$. Blok 4 stránek od \texttt{0x0000B000} jsou virtuální
stránky \texttt{0xB}, \texttt{0xC}, \texttt{0xD}, \texttt{0xE}. Protože je mapování souvislé,
rámce navazují:
\[
\texttt{0xB}\to\texttt{0xA1017},\quad
\texttt{0xC}\to\texttt{0xA1018},\quad
\texttt{0xD}\to\texttt{0xA1019},\quad
\texttt{0xE}\to\texttt{0xA101A}.
\]

Příznaky A a~D: zápis proběhl na \texttt{0x0000D854} (stránka \texttt{0xD}) $\Rightarrow$
stránka \texttt{0xD} je Accessed i~Dirty. Čtení proběhlo ze stránky \texttt{0xC}
$\Rightarrow$ stránka \texttt{0xC} je jen Accessed (A=1, D=0). Stránky \texttt{0xB}
a~\texttt{0xE} nebyly použity (A=0, D=0). Protože jde o~heap, je u~všech P=1, W=1, X=0.

\smallskip
\noindent\textbf{C.1 -- adresy položek.} Tabulka začíná na \texttt{0x13400000}, položka má
4\,B, adresa položky $=\text{báze}+\text{číslo\_stránky}\cdot 4$:
\[
\texttt{0xB}\!\to\!\boxed{\texttt{0x1340002C}},\quad
\texttt{0xC}\!\to\!\boxed{\texttt{0x13400030}},\quad
\texttt{0xD}\!\to\!\boxed{\texttt{0x13400034}},\quad
\texttt{0xE}\!\to\!\boxed{\texttt{0x13400038}}.
\]

\noindent\textbf{C.2 -- hodnoty položek.} Hodnota $=(\text{PFN}\,{\ll}\,12)\,|\,\text{příznaky}$.
Příznakové byty dle zvoleného rozložení: \texttt{0x03} = P+W (bity 0,1); \texttt{0x23} =
P+W+Accessed (přidán bit 5 $=\texttt{0x20}$); \texttt{0x63} = P+W+A+Dirty (navíc bit 6
$=\texttt{0x40}$).
\begin{center}
\begin{tabular}{|c|c|c|c|}
\hline
\textbf{stránka} & \textbf{rámec (PFN)} & \textbf{A, D} & \textbf{položka} \\
\hline
\texttt{0xB} & \texttt{0xA1017} & A=0, D=0 & $\boxed{\texttt{0xA1017003}}$ \\
\hline
\texttt{0xC} & \texttt{0xA1018} & A=1, D=0 & $\boxed{\texttt{0xA1018023}}$ \\
\hline
\texttt{0xD} & \texttt{0xA1019} & A=1, D=1 & $\boxed{\texttt{0xA1019063}}$ \\
\hline
\texttt{0xE} & \texttt{0xA101A} & A=0, D=0 & $\boxed{\texttt{0xA101A003}}$ \\
\hline
\end{tabular}
\end{center}
Přesné hodnoty příznakových bitů závisejí na zvoleném rozložení položky PTE z~části A;
při jiném rozložení polí by se příznakové byty odpovídajícím způsobem změnily (PFN a~adresy
položek zůstávají stejné).

\medskip
\textit{(V~PDF bez náčrtu řešení; náčrt doplněn k~výkladu okruhu, číselné výsledky ověřeny numericky.)}
